\documentclass[12pt]{article}
\usepackage{fullpage}
\usepackage[T1]{fontenc}
\usepackage[utf8]{inputenc}
\usepackage{lmodern}
\usepackage{microtype}
\usepackage{amsmath,amssymb,amsthm}
\usepackage{graphicx}
\usepackage{booktabs}
\usepackage{hyperref}
\usepackage{url}
\usepackage{color}
\usepackage{xcolor}
\usepackage{enumitem}
\usepackage{amsfonts}

\DeclareMathOperator{\vecop}{vec}
\DeclareMathOperator{\diag}{diag}
\DeclareMathOperator{\range}{range}
\newcommand{\bR}{\mathbb{R}}
\newcommand{\tZ}{\tilde{Z}}
\newcommand{\tK}{\tilde{K}}
\newcommand{\KR}{\odot}

\newtheorem{lemma}{Lemma}
\newtheorem{proposition}{Proposition}
\newtheorem{theorem}{Theorem}

\title{Preconditioned CG for the RKHS-Constrained Mode-$k$ CP Subproblem}
\author{}
\date{}

\begin{document}
\maketitle

\section{Setup and notation}

We follow the notation of the problem. The tensor is
$\mathcal{T}\in\bR^{n_1\times\cdots\times n_d}$, $N=\prod_i n_i$,
$n\equiv n_k$, $M=\prod_{i\neq k}n_i$, $K\in\bR^{n\times n}$ is the
symmetric positive semidefinite (psd) RKHS kernel matrix of mode $k$,
\[
  Z=A_d\KR\cdots\KR A_{k+1}\KR A_{k-1}\KR\cdots\KR A_1\in\bR^{M\times r}
\]
is the Khatri--Rao product of the fixed factor matrices, and
$T\in\bR^{n\times M}$ is the mode-$k$ unfolding with missing entries set
to $0$. The unknown is $W\in\bR^{n\times r}$ and the wanted factor is
$A_k=KW$. We use the column-stacking convention for $\vecop$, so that for
any $X\in\bR^{n\times r}$ the vector $\vecop(X)\in\bR^{nr}$ has entries
\begin{equation}\label{eq:vecconv}
  \big(\vecop(X)\big)_{\,i+n(j-1)}=X_{ij},
  \qquad 1\le i\le n,\ 1\le j\le r .
\end{equation}

The data are unaligned: only $q\ll N$ entries of $\mathcal{T}$ are
observed. Each observed entry of the mode-$k$ unfolding corresponds to a
pair $(i_\ell,m_\ell)\in\{1,\dots,n\}\times\{1,\dots,M\}$ with value
$\tau_\ell:=T_{i_\ell,m_\ell}$, for $\ell=1,\dots,q$. The selection
matrix $S\in\bR^{N\times q}$ is a column subset of $I_N$ such that
$S^{\top}\vecop(T)\in\bR^q$ collects the $q$ observed values. With the
convention \eqref{eq:vecconv} applied to the $n\times M$ unfolding $T$,
the $\ell$-th column of $S$ is the standard basis vector
$e_{\,i_\ell+n(m_\ell-1)}\in\bR^N$.

We must solve the symmetric $nr\times nr$ linear system
\begin{equation}\label{eq:system}
  \underbrace{\Big[(Z\otimes K)^{\top}SS^{\top}(Z\otimes K)
      +\lambda\,(I_r\otimes K)\Big]}_{=:A}\,
  \vecop(W)=(I_r\otimes K)\,\vecop(B),
  \qquad B:=TZ\in\bR^{n\times r}.
\end{equation}
Throughout we assume $d<n,r<q\ll N$ and $n\ll M$, and we must avoid every
operation of order $N$ or $M$.

We write $\tK\in\bR^{q\times n}$ for the matrix whose $\ell$-th row is
$K_{i_\ell,:}$ (the $i_\ell$-th row of $K$), and $\tZ\in\bR^{q\times r}$
for the matrix whose $\ell$-th row is $Z_{m_\ell,:}$ (the $m_\ell$-th row
of $Z$). Both have $q$ rows and are formed without touching all of $K$ or
$Z$ (see \S\ref{sec:precompute}). We use $\circledast$ for the Hadamard
(entrywise) product and $\mathbf 1_q\in\bR^q$ for the all-ones vector.

\section{Structure of the system matrix}

\subsection{The two basic Kronecker identities}

\begin{lemma}\label{lem:kron}
For all $X\in\bR^{n\times r}$,
\begin{equation}\label{eq:ZK-action}
  (Z\otimes K)\,\vecop(X)=\vecop\!\big(K X Z^{\top}\big)\in\bR^{N},
\end{equation}
and for all $Y\in\bR^{n\times M}$,
\begin{equation}\label{eq:ZK-adjoint}
  (Z\otimes K)^{\top}\vecop(Y)=\vecop\!\big(K^{\top}Y Z\big)
  =\vecop\!\big(K Y Z\big)\in\bR^{nr},
\end{equation}
the last equality because $K=K^{\top}$.
\end{lemma}

\begin{proof}
Identity \eqref{eq:ZK-action} is the standard rule
$(P\otimes Q)\vecop(X)=\vecop(QXP^{\top})$ for the column-stacking
$\vecop$, with $P=Z$ ($M\times r$) and $Q=K$ ($n\times n$): indeed, with
the convention \eqref{eq:vecconv},
$\big(\vecop(QXP^{\top})\big)_{a+n(b-1)}=\sum_{i,j}Q_{ai}X_{ij}P_{bj}
=\sum_{i,j}(P\otimes Q)_{a+n(b-1),\,i+n(j-1)}X_{ij}$, which is exactly the
$(a+n(b-1))$-th entry of $(P\otimes Q)\vecop(X)$. Transposing gives
$(Z\otimes K)^{\top}=Z^{\top}\otimes K^{\top}=Z^{\top}\otimes K$, and
applying the same rule with $P=Z^{\top}$ ($r\times M$), $Q=K$ yields
\eqref{eq:ZK-adjoint}.
\end{proof}

\subsection{Reduction to the observed entries}

The middle factor $SS^{\top}\in\bR^{N\times N}$ is the diagonal projector
onto the $q$ observed coordinates; multiplicity is impossible because the
columns of $S$ are distinct standard basis vectors. Hence for any
$u\in\bR^N$,
\begin{equation}\label{eq:proj}
  SS^{\top}u=\sum_{\ell=1}^{q}u_{\,i_\ell+n(m_\ell-1)}\,
  e_{\,i_\ell+n(m_\ell-1)} .
\end{equation}

\begin{proposition}[Compressed matrix--vector product]\label{prop:matvec}
For every $W\in\bR^{n\times r}$, set
\begin{equation}\label{eq:matvec-steps}
  F:=\tK W\in\bR^{q\times r},\quad
  y:=(F\circledast\tZ)\,\mathbf 1_r\in\bR^{q},\quad
  G:=\tK^{\top}\big(\diag(y)\,\tZ\big)\in\bR^{n\times r}.
\end{equation}
Then
\begin{equation}\label{eq:matvec-final}
  A\,\vecop(W)=\vecop\!\big(G+\lambda\,KW\big).
\end{equation}
\end{proposition}

\begin{proof}
By \eqref{eq:ZK-action}, $u:=(Z\otimes K)\vecop(W)=\vecop(KWZ^{\top})$, so
the coordinate of $u$ at index $i_\ell+n(m_\ell-1)$ equals
$(KWZ^{\top})_{i_\ell,m_\ell}=K_{i_\ell,:}\,W\,(Z_{m_\ell,:})^{\top}
=\tK_{\ell,:}\,W\,(\tZ_{\ell,:})^{\top}$. With $F=\tK W$ this is
$\sum_{s=1}^r F_{\ell s}\tZ_{\ell s}=\big((F\circledast\tZ)\mathbf
1_r\big)_\ell=y_\ell$. Therefore, by \eqref{eq:proj},
\[
  SS^{\top}u=\sum_{\ell=1}^{q}y_\ell\,e_{\,i_\ell+n(m_\ell-1)}
  =\vecop(R),\qquad
  R\in\bR^{n\times M},\quad R_{i,m}=\!\!\sum_{\ell:(i_\ell,m_\ell)=(i,m)}\!\!y_\ell .
\]
Applying \eqref{eq:ZK-adjoint},
$(Z\otimes K)^{\top}SS^{\top}u=\vecop(KRZ)$, and entrywise
\[
  (KRZ)_{a b}=\sum_{i,m}K_{a i}R_{i m}Z_{m b}
  =\sum_{\ell=1}^q K_{a,i_\ell}\,y_\ell\,Z_{m_\ell,b}
  =\sum_{\ell=1}^q \tK_{\ell a}\,y_\ell\,\tZ_{\ell b}
  =\big(\tK^{\top}\diag(y)\tZ\big)_{ab}=G_{ab}.
\]
Finally $\lambda(I_r\otimes K)\vecop(W)=\lambda\vecop(KW)$ by
\eqref{eq:ZK-action} with $Z=I_r$. Summing the two contributions gives
\eqref{eq:matvec-final}.
\end{proof}

The point of Proposition~\ref{prop:matvec} is that
$A\,\vecop(W)$ is computed from $\tK\,(q\times n)$, $\tZ\,(q\times r)$, and
$K\,(n\times n)$ alone; no object of size $N$ or $M$ ever appears. Because
$A$ is symmetric (it is a sum of a Gram matrix and a symmetric psd term)
and psd, conjugate gradients applies.

\begin{lemma}\label{lem:spd}
$A$ is symmetric psd, and is positive definite (so \eqref{eq:system} has a
unique solution) whenever $\lambda>0$ and $K$ is positive definite.
\end{lemma}
\begin{proof}
For any $v=\vecop(W)$,
$v^{\top}Av=\|S^{\top}(Z\otimes K)v\|_2^2
+\lambda\,v^{\top}(I_r\otimes K)v\ge 0$, since the first term is a squared
norm and $I_r\otimes K\succeq0$ because $K\succeq0$. If $\lambda>0$ and
$K\succ0$ then $I_r\otimes K\succ0$, so $v^{\top}Av\ge\lambda\,
\sigma_{\min}(K)\|v\|_2^2>0$ for $v\ne0$, i.e.\ $A\succ0$.
\end{proof}

\section{The preconditioner}\label{sec:prec}

\subsection{Definition}

We precondition by the \emph{full-data} normal matrix, i.e.\ the matrix
that would arise from \eqref{eq:system} if all $N$ entries were observed,
in which case $SS^{\top}=I_N$. Using
$(Z\otimes K)^{\top}(Z\otimes K)=(Z^{\top}Z)\otimes K^2$,
\begin{equation}\label{eq:Pdef}
  P:=(Z^{\top}Z)\otimes K^{2}+\lambda\,(I_r\otimes K).
\end{equation}
This is an excellent approximation of $A$: $A$ and $P$ share the identical
regularization block $\lambda(I_r\otimes K)$, and their data-fit blocks
are the partial Gram $(Z\otimes K)^{\top}SS^{\top}(Z\otimes K)$ versus the
full Gram $(Z\otimes K)^{\top}(Z\otimes K)$; the two agree up to the
missing-data sampling and, in particular, $A\preceq P$ entrywise in the
Loewner order on the data-fit part, while the regularizer is captured
exactly. Numerically (\S\ref{sec:complexity}) $P^{-1}A$ has a condition
number of order ten in regimes where $A$ has condition number in the
thousands.

\subsection{Cheap factorization of $P$ via two small eigendecompositions}

The matrix $Z^{\top}Z$ is computed \emph{without} $M$-cost: since $Z$ is a
Khatri--Rao product,
\begin{equation}\label{eq:ZtZ}
  Z^{\top}Z=\big(A_d^{\top}A_d\big)\circledast\cdots\circledast
  \big(A_{k+1}^{\top}A_{k+1}\big)\circledast
  \big(A_{k-1}^{\top}A_{k-1}\big)\circledast\cdots\circledast
  \big(A_1^{\top}A_1\big)\in\bR^{r\times r}.
\end{equation}
This is the standard identity that the Gram of a Khatri--Rao product is
the Hadamard product of the factor Grams; each $A_i^{\top}A_i$ costs
$O(n_i r^2)$, so \eqref{eq:ZtZ} costs $O\!\big(r^2\sum_{i\neq k}n_i\big)$,
which is far below $M$.

Now form the symmetric eigendecompositions of the two \emph{small}
matrices,
\begin{equation}\label{eq:eig}
  Z^{\top}Z=U\,D\,U^{\top}\ (D=\diag(d_1,\dots,d_r)\succeq0),\qquad
  K=V\,L\,V^{\top}\ (L=\diag(\ell_1,\dots,\ell_n)\succeq0),
\end{equation}
with $U\in\bR^{r\times r}$, $V\in\bR^{n\times n}$ orthogonal. Because
$U\otimes V$ is orthogonal and diagonalizes every term of \eqref{eq:Pdef}
simultaneously,
\begin{equation}\label{eq:Pdiag}
  P=(U\otimes V)\,
  \Big[\,D\otimes L^{2}+\lambda\,(I_r\otimes L)\,\Big]\,
  (U\otimes V)^{\top},
\end{equation}
where the bracketed matrix is diagonal with entries
\begin{equation}\label{eq:Pevals}
  p_{a,b}:=d_a\,\ell_b^{2}+\lambda\,\ell_b,
  \qquad a=1,\dots,r,\ b=1,\dots,n,
\end{equation}
placed at vec-index $b+n(a-1)$ by \eqref{eq:vecconv}.

\begin{lemma}\label{lem:Pdiag}
Identity \eqref{eq:Pdiag} holds, and \eqref{eq:Pevals} lists the
eigenvalues of $P$.
\end{lemma}
\begin{proof}
$(U\otimes V)^{\top}\big[(Z^{\top}Z)\otimes K^2\big](U\otimes V)
=(U^{\top}Z^{\top}ZU)\otimes(V^{\top}K^2V)=D\otimes L^2$ and
$(U\otimes V)^{\top}\big[I_r\otimes K\big](U\otimes V)
=(U^{\top}U)\otimes(V^{\top}KV)=I_r\otimes L$, using
$V^{\top}K^2V=(V^{\top}KV)^2=L^2$. Summing and conjugating back gives
\eqref{eq:Pdiag}; the diagonal entries of $D\otimes L^2+\lambda(I_r\otimes
L)$ are exactly \eqref{eq:Pevals}, and since $U\otimes V$ is orthogonal
these are the eigenvalues of $P$.
\end{proof}

\subsection{Applying $P^{-1}$ (the preconditioner solve)}

Given a residual $\vecop(R_0)$ with $R_0\in\bR^{n\times r}$, the
preconditioner step $\vecop(Z_0)=P^{-1}\vecop(R_0)$ is computed by
\eqref{eq:Pdiag} as follows. Let $C:=V^{\top}R_0\,U\in\bR^{n\times r}$
(this is the residual expressed in the eigenbasis, since $(U\otimes
V)^{\top}\vecop(R_0)=\vecop(V^{\top}R_0U)$ by \eqref{eq:ZK-action}).
Divide entrywise by the eigenvalues,
\begin{equation}\label{eq:Cdiv}
  \widehat C_{b,a}:=
  \begin{cases} C_{b,a}/p_{a,b}, & p_{a,b}>\varepsilon_P,\\[2pt]
                0, & p_{a,b}\le\varepsilon_P,\end{cases}
  \qquad \varepsilon_P:=\epsilon_{\mathrm{mach}}\max_{a,b}p_{a,b},
\end{equation}
and return $Z_0=V\,\widehat C\,U^{\top}$, i.e.\
$\vecop(Z_0)=(U\otimes V)\vecop(\widehat C)$. The floor in
\eqref{eq:Cdiv} replaces $P^{-1}$ by the Moore--Penrose pseudoinverse on
the (numerical) null space, which occurs precisely when $K$ is rank
deficient ($\ell_b=0\Rightarrow p_{a,b}=0$). This keeps the preconditioner
symmetric positive semidefinite and well defined; on
$\range(I_r\otimes K)\supseteq\range(A)\ni\text{RHS}$ it inverts $P$
exactly, which is all CG requires. When $K\succ0$ all $p_{a,b}>0$ and no flooring occurs.

Each application of $P^{-1}$ costs two $n\times n$-by-$n\times r$ products
($V^{\top}R_0$ and $V\widehat C$) and two $n\times r$-by-$r\times r$
products ($\,\cdot\,U$ and $\,\cdot\,U^{\top}$), i.e.\
$O(n^2r+nr^2)$ flops, with the eigendecompositions \eqref{eq:eig} computed
\emph{once} per mode-$k$ subproblem at cost $O(n^3+r^3)$.

\section{Precomputation without $N$- or $M$-cost}\label{sec:precompute}

\paragraph{Observed-row factors $\tK,\tZ$.}
For $\ell=1,\dots,q$ store the index pair $(i_\ell,m_\ell)$ and the value
$\tau_\ell$ from the sparse data structure. Then $\tK_{\ell,:}=K_{i_\ell,:}$
is read directly from $K$ ($O(qn)$ total) and $\tZ_{\ell,:}=Z_{m_\ell,:}$
is the corresponding row of the Khatri--Rao product, which is computed
from the factor rows by
\begin{equation}\label{eq:Zrow}
  \tZ_{\ell,:}=(A_d)_{j_d,:}\circledast\cdots\circledast
  (A_{k+1})_{j_{k+1},:}\circledast (A_{k-1})_{j_{k-1},:}\circledast\cdots
  \circledast (A_1)_{j_1,:},
\end{equation}
where $(j_1,\dots,j_{k-1},j_{k+1},\dots,j_d)$ is the multi-index of the
column $m_\ell$ of the unfolding (obtained by the mixed-radix
factorization of $m_\ell-1$ with radices $\{n_i\}_{i\neq k}$). Each row
costs $O(dr)$, so $\tZ$ costs $O(qdr)$, well below $M$.

\paragraph{Right-hand side $B=TZ$ (MTTKRP).}
Because $T$ has only the $q$ observed nonzeros $\tau_\ell$ at
$(i_\ell,m_\ell)$,
\begin{equation}\label{eq:B}
  B_{i,:}=\sum_{m}T_{i,m}Z_{m,:}
  =\sum_{\ell:\,i_\ell=i}\tau_\ell\,\tZ_{\ell,:},
\end{equation}
so $B$ is assembled by scatter-adding the $q$ scaled rows $\tau_\ell
\tZ_{\ell,:}$ into row $i_\ell$ of an $n\times r$ accumulator: cost
$O(qr)$. The final right-hand side of \eqref{eq:system} is $\vecop(KB)$,
cost $O(n^2r)$.

\section{Direct symmetric method (baseline)}

We first record the direct dense solver, which forms $A$ and the RHS
explicitly and applies a symmetric (Cholesky/LDL$^{\top}$) factorization.

\begin{center}
\fbox{\parbox{0.95\textwidth}{
\textbf{Algorithm 1: Direct symmetric solver for the mode-$k$ subproblem}\\[2pt]
\textbf{Input:} sparse observations $\{(i_\ell,m_\ell,\tau_\ell)\}_{\ell=1}^q$;
kernel $K\in\bR^{n\times n}$; factors $\{A_i\}_{i\neq k}$; $\lambda>0$.\\
\textbf{Output:} $W\in\bR^{n\times r}$ with $(KW)$ the new mode-$k$ factor.
\begin{enumerate}[leftmargin=2.2em,itemsep=1pt]
  \item Build $\tK\in\bR^{q\times n}$, $\tZ\in\bR^{q\times r}$ by
    \eqref{eq:Zrow} and index lookups.\hfill$O(qn+qdr)$
  \item Form $B=TZ$ by \eqref{eq:B} and set RHS $b=\vecop(KB)$.
    \hfill$O(qr+n^2r)$
  \item Form $A$ explicitly:
    for $\ell=1,\dots,q$ accumulate the rank-structured contribution
    $(\tZ_{\ell,:}\!\otimes\!\tK_{\ell,:})^{\top}
     (\tZ_{\ell,:}\!\otimes\!\tK_{\ell,:})$, then add
    $\lambda(I_r\otimes K)$.\hfill$O(q\,n^2r^2)$
  \item Compute Cholesky/LDL$^{\top}$ factorization $A=R^{\top}R$.
    \hfill$O(n^3r^3)$
  \item Solve $R^{\top}R\,\vecop(W)=b$ by two triangular solves;
    reshape to $W$.\hfill$O(n^2r^2)$
  \item \textbf{return} $W$.
\end{enumerate}
\textbf{Cost:} $O(q\,n^2r^2+n^3r^3)$ time, $O(n^2r^2)$ memory for $A$.}}
\end{center}

This is the $O(n^3r^3)$ solve referenced in the problem, plus the
$O(q\,n^2r^2)$ cost of explicitly forming $A$.

\section{Preconditioned CG method}

\begin{center}
\fbox{\parbox{0.95\textwidth}{
\textbf{Algorithm 2: PCG solver for the mode-$k$ subproblem}\\[2pt]
\textbf{Input:} sparse observations $\{(i_\ell,m_\ell,\tau_\ell)\}_{\ell=1}^q$;
kernel $K\in\bR^{n\times n}$; factors $\{A_i\}_{i\neq k}$; $\lambda>0$;
tolerance $\eta$; max iters $I_{\max}$.\\
\textbf{Output:} $W\in\bR^{n\times r}$ with $(KW)$ the new mode-$k$ factor.\\[3pt]
\emph{Precompute (once per subproblem).}
\begin{enumerate}[leftmargin=2.2em,itemsep=1pt]
  \item Build $\tK\in\bR^{q\times n}$, $\tZ\in\bR^{q\times r}$ by
    \eqref{eq:Zrow}.\hfill$O(qn+qdr)$
  \item Form $B$ by \eqref{eq:B}; set $b=\vecop(KB)$.\hfill$O(qr+n^2r)$
  \item Compute $Z^{\top}Z$ by the Hadamard formula \eqref{eq:ZtZ}.
    \hfill$O\!\big(r^2\sum_{i\neq k}n_i\big)$
  \item Eigendecompose $Z^{\top}Z=UDU^{\top}$ and $K=VLV^{\top}$
    \eqref{eq:eig}; precompute $p_{a,b}=d_a\ell_b^2+\lambda\ell_b$ and the
    floor $\varepsilon_P$.\hfill$O(n^3+r^3)$
\end{enumerate}
\emph{Operators (closures).}
\begin{itemize}[leftmargin=2.2em,itemsep=1pt]
  \item $\mathrm{mul}(W)$: $F\!\leftarrow\!\tK W$;\ \
    $y\!\leftarrow\!(F\circledast\tZ)\mathbf 1_r$;\ \
    $G\!\leftarrow\!\tK^{\top}(\diag(y)\tZ)$;\ \
    \textbf{return} $G+\lambda KW$.\hfill$O(qnr+n^2r)$
  \item $\mathrm{psolve}(R_0)$: $C\!\leftarrow\!V^{\top}R_0U$;\ \
    $\widehat C_{b,a}\!\leftarrow\!C_{b,a}/p_{a,b}$ (or $0$ if
    $p_{a,b}\le\varepsilon_P$);\ \ \textbf{return} $V\widehat C U^{\top}$.
    \hfill$O(n^2r+nr^2)$
\end{itemize}
\emph{PCG iteration} (on the matrix variable $W$; reshape $b$ to
$\bR^{n\times r}$).
\begin{enumerate}[leftmargin=2.2em,itemsep=1pt]
  \item $W\!\leftarrow\!0$;\ \ $R\!\leftarrow\!\mathrm{reshape}(b)$;\ \
    $Q\!\leftarrow\!\mathrm{psolve}(R)$;\ \ $H\!\leftarrow\!Q$;\ \
    $\rho\!\leftarrow\!\langle R,Q\rangle$;\ \ $\beta_0\!\leftarrow\!\|R\|_F$.
  \item \textbf{for} $t=1,\dots,I_{\max}$:
    \begin{enumerate}[leftmargin=1.4em,itemsep=1pt]
      \item $Y\!\leftarrow\!\mathrm{mul}(H)$;\ \
        $\alpha\!\leftarrow\!\rho/\langle H,Y\rangle$.
      \item $W\!\leftarrow\!W+\alpha H$;\ \ $R\!\leftarrow\!R-\alpha Y$.
      \item \textbf{if} $\|R\|_F\le\eta\,\beta_0$: \textbf{break}.
      \item $Q\!\leftarrow\!\mathrm{psolve}(R)$;\ \
        $\rho'\!\leftarrow\!\langle R,Q\rangle$;\ \
        $H\!\leftarrow\!Q+(\rho'/\rho)H$;\ \ $\rho\!\leftarrow\!\rho'$.
    \end{enumerate}
  \item \textbf{return} $W$.
\end{enumerate}
\textbf{Cost:} precompute $O(qn+n^3+r^3)$;
each iteration $O(qnr+n^2r+nr^2)$; total
$O\!\big(qn+n^3+r^3+\kappa(qnr+n^2r)\big)$ where $\kappa$ is the number of
iterations. Memory $O(qn+qr+n^2+r^2)$.}}
\end{center}

Here $\langle X,Y\rangle:=\sum_{ij}X_{ij}Y_{ij}=\vecop(X)^{\top}\vecop(Y)$.
We run CG entirely on the matrix variable $W\in\bR^{n\times r}$; this is
algebraically identical to CG on $\vecop(W)$ because all CG operations are
linear combinations and Euclidean inner products, and $\vecop$ is a linear
isometry $\big(\langle X,Y\rangle=\vecop(X)^{\top}\vecop(Y)\big)$.

\subsection{Correctness}

\begin{theorem}\label{thm:correct}
Algorithm~2 implements preconditioned conjugate gradients for
\eqref{eq:system} with system matrix $A$ and preconditioner $P$ of
\eqref{eq:Pdef}. If $\lambda>0$ and $K\succ0$ then $A\succ0$, $P\succ0$,
and PCG converges to the unique solution of \eqref{eq:system}; the energy
norm error contracts as
\begin{equation}\label{eq:cgrate}
  \|\vecop(W_t)-\vecop(W_\star)\|_A
  \le 2\Big(\tfrac{\sqrt{\kappa_2(P^{-1}A)}-1}
                 {\sqrt{\kappa_2(P^{-1}A)}+1}\Big)^{t}
        \|\vecop(W_0)-\vecop(W_\star)\|_A,
\end{equation}
where $\kappa_2(P^{-1}A)$ is the spectral condition number of $P^{-1}A$.
\end{theorem}

\begin{proof}
The operator $\mathrm{mul}$ returns $\,G+\lambda KW\,$, which by
Proposition~\ref{prop:matvec} equals the reshape of
$A\vecop(W)$; hence $\mathrm{mul}$ is the action of $A$. The operator
$\mathrm{psolve}$ returns $V\widehat CU^{\top}$, which by
Lemma~\ref{lem:Pdiag} and \eqref{eq:Cdiv} equals the reshape of
$P^{-1}\vecop(R_0)$ when $K\succ0$ (no flooring), i.e.\ it is the action of
$P^{-1}$. The right-hand side $b=\vecop(KB)=\vecop(KTZ)=(I_r\otimes
K)\vecop(B)$ matches \eqref{eq:system} (using \eqref{eq:ZK-action} with
$Z=I_r$, and \eqref{eq:B} for $B=TZ$). The loop in Algorithm~2 is exactly
the standard PCG recursion expressed through the linear isometry $\vecop$,
so it coincides iterate-by-iterate with PCG for $A\vecop(W)=b$ with
preconditioner $P$.

By Lemma~\ref{lem:spd}, $A\succ0$ when $\lambda>0,K\succ0$; and $P\succ0$
by \eqref{eq:Pevals} since then every $p_{a,b}=d_a\ell_b^2+\lambda\ell_b\ge
\lambda\ell_b\ge\lambda\,\sigma_{\min}(K)>0$. The convergence bound
\eqref{eq:cgrate} is the classical PCG error estimate for a symmetric
positive definite system with symmetric positive definite preconditioner;
its hypotheses ($A\succ0$, $P\succ0$) are verified, and $P^{-1}A$ is
self-adjoint in the $P$-inner product so $\kappa_2(P^{-1}A)$ is real and
$\ge1$.
\end{proof}

\begin{proposition}[Singular kernel]\label{prop:singular}
If $K$ is psd but singular and $\lambda>0$, then \eqref{eq:system} is
consistent and $A\vecop(W)=b$ determines $KW$ uniquely. With the
pseudoinverse flooring \eqref{eq:Cdiv}, Algorithm~2 returns a solution
$W_\star$ of \eqref{eq:system}, and the recovered factor $A_k=KW_\star$ is
the unique minimizer.
\end{proposition}
\begin{proof}
Write $K=VLV^{\top}$ with $L=\diag(\ell_b)$, and partition the columns of
$V$ into $V_+$ ($\ell_b>0$) and $V_0$ ($\ell_b=0$), so
$\range(K)=\range(V_+)$ and $\ker K=\range(V_0)$. Both the data-fit term
$(Z\otimes K)^{\top}SS^{\top}(Z\otimes K)$ and the regularizer
$\lambda(I_r\otimes K)$ vanish on directions $\vecop(W)$ with
$W$ having columns in $\ker K$, because each contains a left factor $K$
(or $K$ inside $\tK^{\top}\!\cdots$); hence $A$ and $b=\vecop(KB)$ both lie
in $\mathcal R:=\{\vecop(X):X\text{ has columns in }\range K\}$ and $A$ is
invertible \emph{restricted} to $\mathcal R$ (there $A\succeq
\lambda\sigma_{\min}^{+}(K)>0$, where $\sigma_{\min}^{+}$ is the smallest
positive eigenvalue of $K$). Thus \eqref{eq:system} is consistent and
$KW$, equivalently the projection of $\vecop(W)$ onto $\mathcal R$, is
unique. By \eqref{eq:Cdiv}, $\mathrm{psolve}$ realizes $P^{\dagger}$, which
acts as $P^{-1}$ on $\mathcal R$ and as $0$ on $\mathcal R^{\perp}$; since
all CG residuals start in $\mathcal R$ (as $b\in\mathcal R$) and
$A,P^{\dagger}$ preserve $\mathcal R$, the entire iteration stays in
$\mathcal R$ and converges to the unique solution there, whose image
$KW_\star=A_k$ is the claimed minimizer.
\end{proof}

\section{Complexity analysis}\label{sec:complexity}

\paragraph{Per-iteration cost.}
The matrix--vector product $\mathrm{mul}$ costs:
$\tK W$ is $q\times n$ times $n\times r$ at $O(qnr)$;
$(F\circledast\tZ)\mathbf1_r$ is $O(qr)$;
$\tK^{\top}(\diag(y)\tZ)$ is $n\times q$ times $q\times r$ at $O(qnr)$;
$\lambda KW$ is $O(n^2r)$. Total $O(qnr+n^2r)$.
The preconditioner solve $\mathrm{psolve}$ costs $O(n^2r+nr^2)$
(two $n^2r$ and two $nr^2$ products). The CG vector operations are
$O(nr)$. Hence one PCG iteration costs
\begin{equation}\label{eq:periter}
  O\big(qnr+n^2r+nr^2\big).
\end{equation}

\paragraph{Setup cost.}
$\tK$: $O(qn)$; $\tZ$ and $B$: $O(qdr)$; $Z^{\top}Z$ via \eqref{eq:ZtZ}:
$O(r^2\sum_{i\neq k}n_i)\le O(r^2\,(d-1)\,n_{\max})$ where $n_{\max}=\max_i
n_i$ (this is bounded by $O(r^2 dn_{\max})$ and crucially is independent
of $M$); the two eigendecompositions: $O(n^3+r^3)$; RHS $\vecop(KB)$:
$O(n^2r)$. Total setup
\begin{equation}\label{eq:setup}
  O\big(qn+qdr+r^2 d\,n_{\max}+n^3+r^3\big).
\end{equation}

\paragraph{Total.}
With $\kappa$ PCG iterations (which by \eqref{eq:cgrate} is
$O\!\big(\sqrt{\kappa_2(P^{-1}A)}\,\log(1/\eta)\big)$ and is small because
$P$ is a tight preconditioner),
\begin{equation}\label{eq:total}
  \text{PCG total}=
  O\Big(\underbrace{qn+r^2 dn_{\max}+n^3+r^3}_{\text{setup}}
        +\kappa\,(qnr+n^2r+nr^2)\Big).
\end{equation}
Compared with the direct method's $O(q\,n^2r^2+n^3r^3)$
(Algorithm~1), the PCG method removes the dominant $n^3r^3$ factorization
and the $n^2r^2$-per-observation assembly, replacing them by
$O(n^3+r^3)$ setup and $O(\kappa qnr)$ work. \emph{No step touches an
object of size $N=\prod_i n_i$ or $M=\prod_{i\neq k}n_i$}: the only
appearances of the data are through the $q$ observed triples and the small
factor matrices, and $Z^{\top}Z$ is obtained from the Hadamard formula
\eqref{eq:ZtZ}.

\section{Numerical verification}

The following facts were verified numerically against dense reference
computations (random psd $K$, Khatri--Rao $Z$, random observation masks):
(i) the compressed matvec \eqref{eq:matvec-final} reproduces
$A\,\vecop(W)$ to machine precision; (ii) the Hadamard identity
\eqref{eq:ZtZ} and the MTTKRP identity \eqref{eq:B} hold exactly;
(iii) the diagonalization \eqref{eq:Pdiag}--\eqref{eq:Pevals} and the
pseudo-solve \eqref{eq:Cdiv} reproduce $P^{-1}$ (resp.\ $P^{\dagger}$ on
$\range$) to machine precision; (iv) in a representative regime
($n=30$, $r=6$, $q=1500$, $\lambda=10^{-2}$, kernel of rank $20$ plus
ridge), $\kappa_2(A)\approx 1.3\times10^{3}$ while
$\kappa_2(P^{-1}A)\approx 10$, and PCG reached relative residual
$10^{-8}$ in $41$ iterations whereas unpreconditioned CG did not converge
within $200$. These confirm the correctness of Propositions
\ref{prop:matvec}, \ref{prop:singular}, Lemma~\ref{lem:Pdiag}, and the
effectiveness of the preconditioner $P$.

\end{document}
